\section*{Capítulo \ref{ch:g10:cell-metabolism} --- O metabolismo celular}

\begin{solution}{exo:g10:cell-metabolism:1}
O \omterm{def:g10:cell-metabolism:metabolism}{metabolismo} é o conjunto de todas as reações químicas de uma \omterm{def:g10:cells-common-unit:cell}{célula},
que degradam moléculas para obter energia e constroem as moléculas
próprias da \omterm{def:g10:cells-common-unit:cell}{célula}. Uma \omterm{def:g10:cell-metabolism:metabolism}{enzima} é uma \omterm{def:g10:chemistry-of-life:families}{proteína} que acelera uma reação
específica, ligando-se ao substrato dela e liberando o produto, sem ser
consumida.
\end{solution}

\begin{solution}{exo:g10:cell-metabolism:2}
Respiração: $\mathrm{C_6H_{12}O_6} + 6\,\mathrm{O_2} \to
6\,\mathrm{CO_2} + 6\,\mathrm{H_2O}$ + energia. \omterm{def:g10:cell-metabolism:autotrophy}{Fotossíntese}:
$6\,\mathrm{CO_2} + 6\,\mathrm{H_2O}$ + luz $\to \mathrm{C_6H_{12}O_6}
+ 6\,\mathrm{O_2}$. O balanço global de matéria é invertido; mas as
reações, as \omterm{def:g10:cell-metabolism:metabolism}{enzimas}, os lugares (\omterm{def:g10:cells-common-unit:organelle}{mitocôndria} contra \omterm{def:g10:cells-common-unit:organelle}{cloroplasto}) e as
fontes de energia (ligações químicas contra luz) são inteiramente
diferentes.
\end{solution}

\begin{solution}{exo:g10:cell-metabolism:3}
Heterótrofos: cogumelo, \omterm{def:g10:cells-common-unit:cell}{célula} de raiz de cenoura, bactéria do iogurte,
\omterm{def:g10:cells-common-unit:cell}{célula} muscular humana. Autótrofos: musgo, clorela.
\end{solution}

\begin{solution}{exo:g10:cell-metabolism:4}
O gás carbônico infla a bexiga; o etanol se acumula no líquido. Falta
oxigênio ao frasco: com ar borbulhado, a levedura respiraria, não
faria etanol nenhum e usaria muito menos açúcar.
\end{solution}

\begin{solution}{exo:g10:cell-metabolism:5}
A catalase é uma \omterm{def:g10:chemistry-of-life:families}{proteína}; a fervura destrói a forma dela e, portanto,
a atividade dela, e a água oxigenada não se decompõe depressa sozinha.
\end{solution}

\begin{solution}{exo:g10:cell-metabolism:6}
$(7.9 - 0.9)/7 = \qty{1.0}{mg/L}$ por minuto. A curva se achata porque
o oxigênio se esgotou; a levedura passa então à \omterm{def:g10:cell-metabolism:fermentation}{fermentação}, que não
consome oxigênio.
\end{solution}

\begin{solution}{exo:g10:cell-metabolism:7}
$0.8 + 0.2 = \qty{1.0}{mg/L}$ por minuto: a inclinação na luz é a
produção menos a respiração que continua na luz.
\end{solution}

\begin{solution}{exo:g10:cell-metabolism:8}
\qty{18}{g} são \qty{0.1}{mol} de glicose; ela precisa de
\qty{0.6}{mol} de oxigênio, $0.6 \times 32 = \qty{19.2}{g}$, e libera
\qty{0.6}{mol} de gás carbônico, $0.6 \times 44 = \qty{26.4}{g}$.
\end{solution}

\begin{solution}{exo:g10:cell-metabolism:9}
As ervilhas em germinação respiram as reservas delas e liberam calor;
as ervilhas fervidas estão mortas, com as \omterm{def:g10:cell-metabolism:metabolism}{enzimas} destruídas, de modo
que não há reação nem calor. Hermeticamente fechada, a garrafa ficaria
sem oxigênio e a respiração pararia.
\end{solution}

\begin{solution}{exo:g10:cell-metabolism:10}
A metade exposta cora de azul-escuro (amido fabricado pela \omterm{def:g10:cell-metabolism:autotrophy}{fotossíntese}
na luz); a metade coberta não. As 48 horas de escuro deixam a folha
gastar o amido anterior, de modo que qualquer amido encontrado se deve
à luz do dia.
\end{solution}

\begin{solution}{exo:g10:cell-metabolism:11}
A \omterm{def:g10:cell-metabolism:fermentation}{fermentação} rende cerca de quinze vezes menos ATP por glicose que a
respiração; para obter o ATP de que precisa, uma \omterm{def:g10:cells-common-unit:cell}{célula} que fermenta
tem de degradar quinze vezes mais açúcar por minuto.
\end{solution}

\begin{solution}{exo:g10:cell-metabolism:12}
Durante a corrida a necessidade de ATP do músculo superou o oxigênio
que o sangue conseguia entregar; as \omterm{def:g10:cells-common-unit:cell}{células} passaram à \omterm{def:g10:cell-metabolism:fermentation}{fermentação}
láctica, cujo produto, o ácido láctico, se acumulou. Num trote leve o
fornecimento de oxigênio acompanha a demanda, a respiração sozinha
fornece o ATP e nenhum ácido é fabricado.
\end{solution}

\begin{solution}{exo:g10:cell-metabolism:13}
O ambiente mudou (sem luz, com alimento orgânico); os \omterm{def:g10:universal-dna:gene}{genes} não. Os
descendentes reconstroem os \omterm{def:g10:cells-common-unit:organelle}{cloroplastos}, de modo que as instruções
nunca foram perdidas: a perda foi de equipamento em uso, e não das
possibilidades que os \omterm{def:g10:universal-dna:gene}{genes} permitem.
\end{solution}

\begin{solution}{exo:g10:cell-metabolism:14}
A planta faz \omterm{def:g10:cell-metabolism:autotrophy}{fotossíntese} na luz, fornecendo oxigênio e \omterm{def:g10:chemistry-of-life:mineral-organic}{matéria
orgânica} (as folhas que o caramujo come); o caramujo respira,
fornecendo gás carbônico à planta. No escuro a planta apenas respira:
os dois consomem oxigênio, ninguém o produz, e os dois morrem em
poucos dias.
\end{solution}

\begin{solution}{exo:g10:cell-metabolism:15}
$30/2 = 15$ vezes mais glicose por minuto. Os cervejeiros querem
etanol, que só a \omterm{def:g10:cell-metabolism:fermentation}{fermentação} fabrica; uma levedura que respirasse
transformaria o açúcar em gás carbônico, água e mais levedura, e
nenhum álcool.
\end{solution}

\begin{solution}{pb:g10:cell-metabolism:1}
\textbf{1.} A \omterm{def:g10:cell-metabolism:fermentation}{fermentação} alcoólica, porque a cuba fechada não contém
oxigênio: $\mathrm{C_6H_{12}O_6} \to 2\,\mathrm{C_2H_5OH} +
2\,\mathrm{CO_2}$.

\textbf{2.} $200/180 \approx \qty{1.11}{mol}$.

\textbf{3.} Etanol $2 \times 1.11 = \qty{2.22}{mol}$, ou seja $2.22
\times 46 \approx \qty{102}{g}$; volume $102/0.79 \approx
\qty{129}{mL}$; cerca de 13\% em volume.

\textbf{4.} $2.22 \times 44 \approx \qty{98}{g}$ de gás carbônico;
$2.22 \times 24 \approx \qty{53}{L}$ de gás.

\textbf{5.} Cerca de \qty{53000}{L} $= \qty{53}{m^3}$. A adega tem
$4 \times 6 \times 2.5 = \qty{60}{m^3}$: o gás liberado quase a
encheria.

\textbf{6.} $\mathrm{C_6H_{12}O_6} + 6\,\mathrm{O_2} \to
6\,\mathrm{CO_2} + 6\,\mathrm{H_2O}$.

\textbf{7.} Oxigênio $6 \times 1.11 = \qty{6.67}{mol}$, $6.67 \times 32
\approx \qty{213}{g}$; gás carbônico $6.67 \times 44 \approx
\qty{293}{g}$.

\textbf{8.} Três vezes mais: a respiração transforma os seis carbonos
da glicose em gás carbônico, a \omterm{def:g10:cell-metabolism:fermentation}{fermentação} apenas dois, e os outros
quatro ficam no etanol.

\textbf{9.} Com a respiração, quinze vezes mais tempo: cada glicose dá
quinze vezes mais ATP.

\textbf{10.} O etanol só é fabricado pela \omterm{def:g10:cell-metabolism:fermentation}{fermentação}; uma cuba aerada
daria gás carbônico, água e levedura, e nenhum vinho.

\textbf{11.} $1.11 \times 70 \approx \qty{78}{kJ}$ por litro;
\qty{78000}{kJ}, ou seja \qty{78}{MJ}, para a cuba.

\textbf{12.} $78/4.2 \approx \qty{18.5}{\celsius}$.

\textbf{13.} $20 + 18.5 \approx \qty{38.5}{\celsius}$: no limite da
levedura, e uma adega quente empurraria a cuba para além dele. Uma
cuba grande perde calor devagar porque a superfície dela é pequena
comparada com o volume, de modo que ela precisa ser refrigerada de
propósito.

\textbf{14.} As \omterm{def:g10:cell-metabolism:metabolism}{enzimas} são \omterm{def:g10:chemistry-of-life:families}{proteínas}; o calor as desdobra e elas
param de funcionar. Sem as \omterm{def:g10:cell-metabolism:metabolism}{enzimas} dela, a levedura não consegue
transformar o açúcar, por mais que ainda reste.

\textbf{15.} Nas ligações químicas do etanol: cerca de \qty{2800}{kJ}
por mol de glicose, e é por isso que o etanol queima.

\textbf{16.} No chão, numa camada que engrossa à medida que a
\omterm{def:g10:cell-metabolism:fermentation}{fermentação} avança. Uma vela levada rente ao chão se apaga onde o gás
carbônico deslocou o oxigênio --- um aviso antes que uma pessoa seja
afetada.

\textbf{17.} A \omterm{def:g10:cell-metabolism:fermentation}{fermentação} láctica. Ela rende quinze vezes menos ATP; o
cérebro e o coração não conseguem suprir as necessidades deles desse
jeito, e o ácido láctico se acumula. A perda de consciência vem em um
ou dois minutos.

\textbf{18.} O etanol em concentração alta danifica as membranas e as
\omterm{def:g10:cell-metabolism:metabolism}{enzimas} da levedura: o próprio produto das \omterm{def:g10:cells-common-unit:cell}{células} envenena o
\omterm{def:g10:cell-metabolism:metabolism}{metabolismo} delas.

\textbf{19.} As bolhas são de gás carbônico vindo da \omterm{def:g10:cell-metabolism:fermentation}{fermentação}. O
etanol evapora no forno (ele ferve a \qty{78}{\celsius}).

\textbf{20.} Cerca de 13\% de álcool em volume e \qty{53}{L} de gás por
litro de mosto; os dois decorrem de a levedura fermentar, sem oxigênio,
cada glicose em dois etanóis e dois gases carbônicos.
\end{solution}
